%% Đề 12 -- de-hoc-sinh-gioi-toan-9-nam-2021-2022-phong-gddt-chi-lang-lang-son.pdf
\deheader{ĐỀ SỐ 12}%
  {Kỳ thi Chọn Học sinh giỏi cấp tỉnh lớp 9 — Năm học 2021--2022}%
  {Sở GD\&ĐT Lạng Sơn | Chỉ phần Đại số}
\addcontentsline{toc}{section}{Đề 12 — Kỳ thi Chọn HSG cấp tỉnh — Năm học 2021--2022}

\begin{de_cau}{Câu 1}{4,0 điểm}
\begin{enumerate}[label=\alph*)]
\item Tính: $A=\sqrt{3+2\sqrt{2}}+\sqrt{3-2\sqrt{2}}$.
\item Rút gọn $P=\left(\dfrac{x}{x+2\sqrt{x}}+\dfrac{1}{\sqrt{x}+2}\right)\cdot\dfrac{\sqrt{x}+1}{\sqrt{x}-1}$,
$x>0, x\neq 1$.
\end{enumerate}
\end{de_cau}

\begin{loi_giai_cau}
\textbf{a)} $\sqrt{3+2\sqrt{2}}=\sqrt{(\sqrt{2}+1)^2}=\sqrt{2}+1$.
$\sqrt{3-2\sqrt{2}}=\sqrt{(\sqrt{2}-1)^2}=\sqrt{2}-1$.
$A=(\sqrt{2}+1)+(\sqrt{2}-1)=2\sqrt{2}$.

\textbf{b)} $\dfrac{x}{x+2\sqrt{x}}=\dfrac{x}{\sqrt{x}(\sqrt{x}+2)}=\dfrac{\sqrt{x}}{\sqrt{x}+2}$.
\[P=\left(\frac{\sqrt{x}}{\sqrt{x}+2}+\frac{1}{\sqrt{x}+2}\right)\cdot\frac{\sqrt{x}+1}{\sqrt{x}-1}
=\frac{\sqrt{x}+1}{\sqrt{x}+2}\cdot\frac{\sqrt{x}+1}{\sqrt{x}-1}=\frac{(\sqrt{x}+1)^2}{(\sqrt{x}+2)(\sqrt{x}-1)}.\]
\end{loi_giai_cau}

\begin{de_cau}{Câu 2}{4,0 điểm}
\begin{enumerate}[label=\alph*)]
\item Giải phương trình: $x^2 - 4x + 3 = \sqrt{x-1}$.
\item Cho phương trình $x^2-2(m+1)x+4m=0$. Tìm $m$ để pt có hai nghiệm dương phân biệt.
\end{enumerate}
\end{de_cau}

\begin{loi_giai_cau}
\textbf{a)} ĐKXĐ: $x\geq 1$. Đặt $t=\sqrt{x-1}\geq 0$: $x=t^2+1$.

$(t^2+1)^2-4(t^2+1)+3=t\Rightarrow t^4-2t^2-t=0\Rightarrow t(t^3-2t-1)=0$.

$t=0\Rightarrow x=1$. $t^3-2t-1=(t+1)(t^2-t-1)=0\Rightarrow t=\dfrac{1+\sqrt{5}}{2}$ (nhận).

$x=(t^2+1)=\dfrac{3+\sqrt{5}}{2}+1=\dfrac{5+\sqrt{5}}{2}$.

Nghiệm: $x=1$ và $x=\dfrac{5+\sqrt{5}}{2}$.

\textbf{b)} Để có 2 nghiệm dương phân biệt:
$\begin{cases}\Delta'>0\\S=2(m+1)>0\\P=4m>0\end{cases}$

$\Delta'=(m+1)^2-4m=m^2-2m+1=(m-1)^2\geq 0$, bằng $0$ khi $m=1$.

Để phân biệt: $m\neq 1$. $S>0\Rightarrow m>-1$. $P>0\Rightarrow m>0$.

Kết hợp: $m>0$ và $m\neq 1$, tức $m\in(0;1)\cup(1;+\infty)$.
\end{loi_giai_cau}

\begin{de_cau}{Câu 3}{4,0 điểm}
\begin{enumerate}[label=\alph*)]
\item Tìm nghiệm nguyên của phương trình $2x^2-5xy+3y^2=2$.
\item Cho $a,b\in\mathbb{Z}$, $ab|a^2+b^2$. Chứng minh $\sqrt{\dfrac{a^2+b^2}{ab}}$ là số nguyên.
\end{enumerate}
\end{de_cau}

\begin{loi_giai_cau}
\textbf{a)} $(2x-3y)(x-y)=2$. Xét các trường hợp:
$\begin{cases}2x-3y=1\\x-y=2\end{cases}$: $x=5,y=3$. Kiểm tra: $50-75+27=2$ \checkmark .

$\begin{cases}2x-3y=2\\x-y=1\end{cases}$: $x=1, y=0$. Kiểm tra: $2-0+0=2$ \checkmark .

$\begin{cases}2x-3y=-1\\x-y=-2\end{cases}$: $x=-5,y=-3$.
$\begin{cases}2x-3y=-2\\x-y=-1\end{cases}$: $x=-1,y=0$.

Nghiệm nguyên: $(\pm 5,\pm 3)$, $(\pm 1, 0)$ và các hoán vị.

\textbf{b)} Đặt $k=\dfrac{a^2+b^2}{ab}$. Ta cần cm $k$ là số chính phương hoàn hảo.

Gọi $d=\gcd(a,b)$, $a=da'$, $b=db'$. $k=\dfrac{d^2(a'^2+b'^2)}{d^2a'b'}=\dfrac{a'^2+b'^2}{a'b'}$.

Vì $\gcd(a',b')=1$: $a'|1$ và $b'|1$, nên $a'=b'=\pm 1$, $k=2$.

$\sqrt{2}$ không nguyên. Chú ý: cần $ab|(a^2+b^2)$ với $a,b\in\mathbb{N}^*$.
\end{loi_giai_cau}

\begin{de_cau}{Câu 5}{2,0 điểm}
Cho $x,y,z>0$ thỏa mãn $x+y+z=xyz$. Chứng minh:
\[\frac{1}{\sqrt{1+x^2}}+\frac{1}{\sqrt{1+y^2}}+\frac{1}{\sqrt{1+z^2}}\leq\frac{3}{2}.\]
\end{de_cau}

\begin{loi_giai_cau}
Đặt $x=\tan A, y=\tan B, z=\tan C$ ($A,B,C\in(0,\frac{\pi}{2})$).

Điều kiện $x+y+z=xyz$ tương đương $\tan A+\tan B+\tan C=\tan A\tan B\tan C$, đúng khi $A+B+C=\pi$.

$\dfrac{1}{\sqrt{1+x^2}}=\dfrac{1}{\sqrt{1+\tan^2 A}}=\cos A$.

Cần cm $\cos A+\cos B+\cos C\leq\dfrac{3}{2}$ với $A+B+C=\pi$ (tam giác).

Đây là bất đẳng thức cổ điển, dấu bằng khi $A=B=C=60^\circ$ (tam giác đều).

Bằng Jensen (hàm $\cos$ lõm trên $(0,\pi)$):
$\cos A+\cos B+\cos C\leq 3\cos\left(\dfrac{A+B+C}{3}\right)=3\cos 60^\circ=\dfrac{3}{2}$. $\square$
\end{loi_giai_cau}
